
\subsubsection{3.1 Problem Formulation}

Let $D = \{(q_i, r_i^+, r_i^-, \ell_i, t_i)\}$ denote a preference dataset where $q_i$ is a prompt, $r_i^+$ and $r_i^-$ are preferred and rejected responses, $\ell_i$ is the preference label, and $t_i \in \{1, 2, 3\}$ is the annotation stratum. The preference model is:

$$P(\ell_i = 1 \mid x_i, t_i) = \sigma\!\left(\beta_Q^{(t)} \cdot Q_{\text{early}}(x_i) + \beta_L^{(t)} \cdot \Delta L_i + \beta_H^{(t)} \cdot \Delta H_i + \beta_S^{(t)} \cdot \Delta S_i\right)$$

The \textbf{AAI drift signal} is $(\Delta\beta_L, \Delta\beta_H, \Delta\beta_S)$ where $\Delta\beta_k = \beta_k^{(3)} - \beta_k^{(1)}$. A coefficient is classified as directionally drifted when the 95\% bootstrap confidence intervals for strata 1 and 3 are non-overlapping.

\subsubsection{3.2 Quality Covariate: Q_early}

Q_early is a logistic regression model trained exclusively on stratum-1 preference labels using non-stylistic features and then frozen. Its prediction is included as a covariate in all downstream regressions to decompose stylistic preference change from quality recalibration. Stability criterion: |β_Q| < 0.2. In the H-M2 experiment, the observed β_Q in the late stratum is −0.017 (within threshold). Q_early is calibrated via Platt scaling.

\subsubsection{3.3 Stylistic Feature Extraction}

\begin{table}[htbp]
\centering
\begin{tabular}{lll}
\toprule
Feature & Symbol & Operationalization \\
\midrule
Verbosity & Δ_L & word_count(r⁺) − word_count(r⁻) \\
Hedging & Δ_H & hedge_count(r⁺) − hedge_count(r⁻) \\
Structured reasoning & Δ_S & struct_count(r⁺) − struct_count(r⁻) \\
\bottomrule
\end{tabular}
\end{table}

All features are standardized with a shared StandardScaler fit on stratum-1 data only. Variance inflation factors (VIF) are below 1.03 for all features, indicating no multicollinearity.

\subsubsection{3.4 Round-Stratified Coefficient Comparison (H-M2 Protocol)}

For strata t ∈ {1, 3}, logistic regression is fit separately and coefficients are extracted. Bootstrap confidence intervals are computed from 2,000 stratified resamples. The pre-registered gate requires n_directional ≥ 2 of 3 stylistic features showing non-overlapping 95\% CIs.

\subsubsection{3.5 AI-Typicality Geometric Projection (H-M1 Protocol)}

\textbf{AI-typicality vector:} The centroid difference between AI-generated and human-written responses in stratum-1 HH-RLHF, encoded via the frozen all-MiniLM-L6-v2 sentence transformer (384 dimensions), L2-normalized.

\textbf{Projection score:} Cosine projection of each sample's preference gradient onto the AI-typicality vector.

\textbf{Between-group design:} Due to the absence of worker identity metadata in the public WebGPT JSONL release, the planned within-annotator PanelOLS regression was replaced by a between-group tercile comparison. Groups are assigned by |score_0 − score_1| tercile as a proxy for annotator confidence. This fallback was documented in the experiment brief (02c_experiment_brief.md) prior to execution.

\textbf{Discriminant validity:} A placebo permutation test using 200 random unit vectors generates a null distribution for the projection coefficient. The empirical p-value is the fraction of placebo runs yielding a coefficient at least as large as the observed value.

\subsubsection{3.6 Datasets}

\begin{itemize}
\item \textbf{Anthropic HH-RLHF} [Bai et al. 2022]: 160,800 preference comparisons partitioned into three strata of 53,600 rows each by row index. \textbf{Note:} This is index-based partitioning, not genuine temporal metadata; the dataset does not contain per-annotation timestamps or annotator identifiers. The strata are treated as proxies for annotation phase ordering.
\end{itemize}

\begin{itemize}
\item \textbf{OpenAI WebGPT} [Stiennon et al. 2020]: 19,578 preference comparisons loaded from a local JSONL file. Worker identity fields are absent from the public release, preventing within-annotator analysis. A between-group tercile design is used as a documented fallback.
\end{itemize}

